\documentclass[12pt,a4paper]{ctexart}

\usepackage[a4paper,margin=2.2cm]{geometry}
\usepackage{array}
\usepackage{booktabs}
\usepackage{longtable}
\usepackage{tabularx}
\usepackage{multirow}
\usepackage{enumitem}
\usepackage{xcolor}
\usepackage{hyperref}
\usepackage{float}
\usepackage{graphicx}
\usepackage{amsmath}
\usepackage{amssymb}
\usepackage{fancyhdr}
\usepackage{titlesec}
\usepackage{setspace}
\usepackage{ragged2e}

\hypersetup{
  colorlinks=true,
  linkcolor=blue!60!black,
  urlcolor=blue!70!black,
  citecolor=blue!60!black,
  pdftitle={MiroFish / CogSec 项目总结报告},
  pdfauthor={Codex},
  pdfsubject={信安竞赛项目总结、系统架构、对标分析与创新性说明},
  pdfkeywords={MiroFish, CogSec, 认知安全, 反事实推演, 个人决策图谱, 社会工程学, 项目总结}
}

\setstretch{1.25}
\pagestyle{fancy}
\fancyhf{}
\lhead{MiroFish / CogSec 项目总结报告}
\rhead{\thepage}
\titleformat{\section}{\Large\bfseries}{\thesection}{0.8em}{}
\titleformat{\subsection}{\large\bfseries}{\thesubsection}{0.8em}{}
\titleformat{\subsubsection}{\normalsize\bfseries}{\thesubsubsection}{0.8em}{}

\newcolumntype{Y}{>{\RaggedRight\arraybackslash}X}
\newcommand{\ProjectName}{MiroFish / CogSec}
\newcommand{\CorePosition}{\textbf{基于反事实推演的个人决策图谱}}
\newcommand{\ResearchDate}{2026-04-21}
\newcommand{\StatusDone}{\textcolor{green!45!black}{[已确认]}}
\newcommand{\StatusPlan}{\textcolor{orange!80!black}{[规划中]}}
\newcommand{\StatusRisk}{\textcolor{red!70!black}{[需补强]}}
\newcommand{\SrcMiroFishRepo}{\footnote{MiroFish GitHub 仓库：666ghj, \emph{MiroFish: A Simple and Universal Swarm Intelligence Engine, Predicting Anything}, \url{https://github.com/666ghj/MiroFish}；MiroFish-Offline GitHub 仓库：nikmcfly, \emph{MiroFish-Offline}, \url{https://github.com/nikmcfly/MiroFish-Offline}。访问日期：\ResearchDate。}}
\newcommand{\SrcCompetitionLandscape}{\footnote{相关公开方案包括：Bitdefender Scamio 官方页 \url{https://www.bitdefender.com/en-us/consumer/scamio}；Norton AI Scam Protection / Genie 官方页 \url{https://hk-en.norton.com/feature/ai-scam-protection}；SoSafe 官方页 \url{https://sosafe-awareness.com/products/phishing-simulations/}；KnowBe4 AIDA 官方页 \url{https://www.knowbe4.com/products/aida}；Adaptive Security 官方页 \url{https://www.adaptivesecurity.com/}；Microsoft Defender Phishing Triage Agent 文档 \url{https://learn.microsoft.com/en-us/defender-xdr/phishing-triage-agent}；Pretexta GitHub 仓库 \url{https://github.com/dalpan/Pretexta}；AXIS 论文页 \url{https://arxiv.org/abs/2505.17801}。访问日期：\ResearchDate。}}
\newcommand{\SrcDualSystem}{\footnote{Xu, F.; Liu, A.; Li, X. \emph{Victimization mechanisms and countermeasures in telecom network fraud: a dual-system theoretical perspective}. Frontiers in Psychology, 2025. \url{https://www.frontiersin.org/journals/psychology/articles/10.3389/fpsyg.2025.1637935/full}。访问日期：\ResearchDate。}}
\newcommand{\SrcUserContext}{\footnote{Dawkins, S.; Jacobs, J. \emph{Phishing for User Context: Understanding the NIST Phish Scale}. NIST / FISSEA, 2023. \url{https://www.nist.gov/publications/phishing-user-context-understanding-nist-phish-scale}；Frauenstein, E. D.; Flowerday, S.; Wiltermuth, M. T. \emph{Human cognition through the lens of social engineering cyberattacks}. Frontiers in Psychology, 2020. \url{https://pmc.ncbi.nlm.nih.gov/articles/PMC7554349/}。访问日期：\ResearchDate。}}
\newcommand{\SrcPersuasion}{\footnote{Pretexta GitHub 仓库：\url{https://github.com/dalpan/Pretexta}；Khadka, K. et al. \emph{A Survey on the Principles of Persuasion as a Social Engineering Strategy in Phishing}. arXiv:2412.18488, \url{https://doi.org/10.48550/arXiv.2412.18488}；Lawson, P. et al. \emph{Interaction of Personality and Persuasion Tactics in Email Phishing Attacks}. DOI: \url{https://doi.org/10.1177/1541931213601815}；Zielinska, O. A. et al. \emph{A Temporal Analysis of Persuasion Principles in Phishing Emails}. DOI: \url{https://doi.org/10.1177/1541931213601175}。访问日期：\ResearchDate。}}
\newcommand{\SrcAxis}{\footnote{Gyevnár, B.; Lucas, C. G.; Albrecht, S. V.; Cohen, S. B. \emph{Integrating Counterfactual Simulations with Language Models for Explaining Multi-Agent Behaviour}. arXiv:2505.17801, \url{https://arxiv.org/abs/2505.17801}。访问日期：\ResearchDate。}}
\newcommand{\SrcConsumerTools}{\footnote{Bitdefender Scamio 官方页：\url{https://www.bitdefender.com/en-us/consumer/scamio}；Norton AI Scam Protection / Genie 官方页：\url{https://hk-en.norton.com/feature/ai-scam-protection}。访问日期：\ResearchDate。}}
\newcommand{\SrcEnterprisePlatforms}{\footnote{SoSafe 官方页：\url{https://sosafe-awareness.com/products/phishing-simulations/}；KnowBe4 AIDA 官方页：\url{https://www.knowbe4.com/products/aida}；Adaptive Security 官方页：\url{https://www.adaptivesecurity.com/}。访问日期：\ResearchDate。}}
\newcommand{\SrcPhishingTriage}{\footnote{Microsoft Learn, \emph{Security Copilot Phishing Triage Agent in Microsoft Defender}, \url{https://learn.microsoft.com/en-us/defender-xdr/phishing-triage-agent}。访问日期：\ResearchDate。}}
\newcommand{\SrcPretexta}{\footnote{dalpan, \emph{Pretexta: Open-source social engineering simulation lab}, GitHub repository, \url{https://github.com/dalpan/Pretexta}。访问日期：\ResearchDate。}}
\newcommand{\SrcDigitalDeception}{\footnote{Alqubaisi, S. et al. \emph{Digital deception: generative artificial intelligence in social engineering and phishing}. Artificial Intelligence Review, 2025. \url{https://link.springer.com/article/10.1007/s10462-024-10973-2}。访问日期：\ResearchDate。}}
\newcommand{\SrcGovernance}{\footnote{Microsoft Defender Phishing Triage Agent 文档：\url{https://learn.microsoft.com/en-us/defender-xdr/phishing-triage-agent}；Pretexta GitHub 仓库：\url{https://github.com/dalpan/Pretexta}。访问日期：\ResearchDate。}}

\begin{document}

\begin{titlepage}
  \centering
  \vspace*{2.2cm}
  {\Huge\bfseries MiroFish / CogSec 项目总结报告\par}
  \vspace{0.6cm}
  {\Large 面向信安竞赛的认知安全系统说明文档\par}
  \vspace{0.5cm}
  {\Large 主题定位：\CorePosition{}\par}
  \vspace{1.8cm}
  {\large 文档版本：v3.2\par}
  \vspace{0.3cm}
  {\large 日期：\ResearchDate\par}
  \vfill
  \begin{minipage}{0.88\textwidth}
  \small
  \textbf{文档摘要：}本文围绕 \ProjectName{} 项目，首先说明为什么可以基于 MiroFish 底座构建一个以反事实推演为核心的个人决策图谱系统；随后对市场同类产品和 GitHub 开源方案进行调研；再结合认知心理学、社会工程学和解释型 AI 的研究基础，说明系统为何采用当前设计；最后对系统架构、工程实现、风险治理、创新点与迭代路线进行完整拆解，形成一份可独立阅读、可直接用于竞赛汇报与书面提交的正式报告。简而言之，该系统面向诈骗、钓鱼和社会工程场景，在用户即将执行高风险动作前，帮助其看清当前风险、比较不同决策路径，并给出可执行的干预建议。
  \end{minipage}
\end{titlepage}

\tableofcontents
\newpage

\section{执行摘要}

\ProjectName{} 当前的核心目标，是将 MiroFish 在复杂情境建模、多智能体推演与结构化报告生成方面的底座能力，收束到一个更聚焦也更适合信安竞赛的方向上：\CorePosition{}。这一定位并不是对原有底座能力的否定，而是将原本用于“世界如何演化”的建模能力，迁移到“个体在高风险场景中会如何做出决策”这一更具体、更可验证的问题上。

\subsection{它是做什么的}

如果用最直接的话来概括，\ProjectName{} 是一个面向诈骗、钓鱼和社会工程场景的\textbf{决策辅助系统}。它不是只判断“这条消息像不像诈骗”，而是进一步判断“用户如果继续按对方要求操作，接下来可能会发生什么；如果此刻停下来核验，结果又会怎样”。系统的目标是在损失真正发生前，把原本模糊的风险变成一张可解释、可比较、可干预的决策图谱。

从使用流程上看，系统主要完成三件事：
\begin{enumerate}[leftmargin=2em]
  \item 接收可疑文本、聊天记录、截图或场景描述，并识别其中的高危动作与诱导线索；
  \item 结合个体风险因子、威胁知识和当前上下文，构造继续执行与暂停核验两条反事实路径；
  \item 输出风险差异、关键转折点、可逆性窗口和个性化建议，帮助用户在执行危险动作前及时调整决策。
\end{enumerate}

\subsection{本地仓库已可直接确认的实现事实}

在真正扩写报告前，对本地仓库的核对表明，当前项目已经具备一条可以直接落笔的主链，而不仅是概念方案。后端 \texttt{backend/app/modules/cognitive\_profiler.py} 中明确给出了 18 维认知画像，并将其划分为 \texttt{state}、\texttt{vulnerability} 和 \texttt{protection} 三类；\texttt{CognitiveProfile} 同时提供了 \texttt{overall\_vulnerability\_score()}、\texttt{summary()} 和 \texttt{to\_persona\_state\_vector()}，说明画像并不是停留在展示层，而是会被压缩成运行时数字孪生向量。

运行时层面，\texttt{runtime\_schema.py} 中已定义 \texttt{PersonaStateVector}、\texttt{RiskGraphBundle}、\texttt{WorldStateSnapshot}、\texttt{ForkComparisonResult} 和 \texttt{InterventionPrescription} 等主对象。也就是说，项目已经在数据结构上把“数字孪生—风险图谱—世界状态—分支比较—干预处方”串成了一条完整链。再往前走，\texttt{ThreatKnowledgeRAG} 不只是检索案例，而是直接负责构造 \texttt{persuasion\_principles}、\texttt{persona\_weakness\_hits}、\texttt{asset\_targets}、\texttt{environment\_context}、\texttt{fork\_points}、\texttt{consistency\_score} 和 \texttt{hallucination\_rollback} 在内的 \texttt{RiskGraphBundle}。这表明“攻击-人-环境图谱”在项目中已经有明确对象承载。

主判别运行时同样不是纸面设想。\texttt{mainline\_runtime.py} 中的 \texttt{MiroFishRuntime} 已实现 \texttt{initialize\_world\_state()}、\texttt{select\_primary\_fork()} 和 \texttt{run()}，并把主分叉规则具体写成 \texttt{transfer\_money}、\texttt{screen\_share}、\texttt{verification\_code}、\texttt{unknown\_app\_download}、\texttt{social\_isolation} 和 \texttt{fake\_official\_verification} 六类。前端 \texttt{CogSecWorkbench.vue} 也已经对应组织成“用户数字孪生—攻击-人-环境图谱—双分支反事实推演—风险 / 可逆性曲线—个性化干预处方”五屏工作台。换言之，后文关于 18 维特征、主链对象、Fork 节点和处方输出的叙述，都可以直接落回本地实现，而不是只来自参考材料。

之所以能够基于 MiroFish 形成这一新方向，原因在于二者在结构上具有天然衔接性。MiroFish 原本擅长处理的是“复杂输入、图谱构建、情境生成、状态演化和报告输出”的闭环\SrcMiroFishRepo。 而诈骗与社会工程学防御虽然看上去更贴近安全场景，但其真正难点同样不是单点检测，而是围绕人物、环境、诱导策略、状态变化和行为后果所形成的复杂决策链。因此，把 MiroFish 的环境建模、状态更新和交互式推演能力收束到“个人关键动作前的风险分支比较”，在方法上是连贯的，在工程上也是可行的。

从公开市场和开源方案来看，目前已存在大量面向消费者的 AI 反诈助手、面向企业的人因风险训练平台、面向 SOC 的 triage agent，以及面向训练或研究的社工仿真与反事实解释系统\SrcCompetitionLandscape。 这些方案在内容识别、员工训练、组织级分析和解释建模方面都已经取得了成熟进展，但较少有系统把“实时个人反诈骗决策图谱 + 关键动作前的双分支反事实推演 + 个体风险因子命中链 + 可逆性窗口”作为核心产品定义。由此可以判断，本项目仍然具备清晰的差异化空间。

在设计依据上，项目并非单纯堆叠模块。双系统理论解释了为什么人在高压、紧迫和不确定条件下更容易转入快速、直觉和启发式的处理模式\SrcDualSystem。 用户情境研究说明了同样的内容对不同个体、不同上下文具有不同风险\SrcUserContext。 社会工程研究和相关开源实践则表明，权威、紧迫、社会认同、承诺一致性、互惠与喜好等说服原则，往往是攻击得以推进的关键触发机制\SrcPersuasion。 反事实解释研究进一步证明，把复杂行为系统的结果改写为“如果换一条路径会怎样”的比较式解释，是一种具有明确研究依据的解释策略\SrcAxis。 这些理论共同支持了项目当前的系统结构：T0 快速护栏负责拦截高危红旗，T1 深度推演负责解释情境、比较路径、输出干预。

从工程实现看，项目已经具备前后端分离骨架、主后端入口与最小 Demo 入口、五屏工作台界面、风险图谱和反事实报告对象，以及针对多个核心模块的测试和 benchmark 钩子。这说明项目并不是停留在概念层面，而已经拥有真实的实现基础。下一阶段最重要的工作，不是继续扩大叙事，而是将“个人决策 DAG、Branch A / Branch B 对照、可逆性曲线、干预窗口与处方报告”这一主线进一步做透并量化验证。

\section{为什么基于 MiroFish 可以构建个人决策图谱}

\subsection{MiroFish 原始能力的结构价值}

GitHub 原版 MiroFish 将自己定义为一种面向未来场景排演的多智能体预测引擎，强调用户输入现实种子材料后，系统可以完成图谱构建、环境搭建、仿真执行、报告生成和深度交互\SrcMiroFishRepo。 从这一点看，MiroFish 的核心并不是某个特定场景，而是一种更一般的\textbf{复杂情境建模与推演框架}。

这一框架至少包含五个对当前项目极其重要的通用能力：
\begin{enumerate}[leftmargin=2em]
  \item 将非结构化输入转化为可计算的图谱或状态对象；
  \item 在多角色、多因素约束下构造情境环境；
  \item 跟踪状态演化与关键分叉；
  \item 将复杂过程转译为可阅读的报告；
  \item 支持进一步交互与解释。
\end{enumerate}

诈骗与社会工程学场景虽然对象更聚焦，但本质上同样需要这五类能力。系统仍要接收复杂输入，仍要描述角色与环境关系，仍要识别状态变化与关键节点，仍要解释结果并给出建议。因此，从底座能力上看，MiroFish 与当前项目之间不是“嫁接关系”，而是“问题域收束关系”。

\subsection{为什么可以从“世界排演”收束到“个体决策”}

原版 MiroFish 更偏重群体、环境和未来世界的演化；而当前项目更偏重个体、关键动作和即时决策。两者表面上看尺度不同，但方法上仍然连贯。因为无论分析对象是一个社会系统还是一个个体用户，核心都涉及：
\begin{enumerate}[leftmargin=2em]
  \item 当前处境如何被结构化表示；
  \item 哪些因素会推动状态变化；
  \item 在什么节点上会出现路径分叉；
  \item 分叉之后的结果将如何不同；
  \item 如何把复杂演化压缩为可解释结论。
\end{enumerate}

因此，从 MiroFish 走向个人决策图谱，并不是抛弃原先的情境建模，而是把“社会环境的平行世界”缩放成“个体决策的局部平行路径”。这种缩放反而更适合竞赛场景，因为它使系统价值更清楚、可验证性更强、展示链路更短。

\subsection{当前项目从 MiroFish 继承了什么}

从工程痕迹上看，当前项目与 MiroFish 存在明确承继关系。依赖中存在 \texttt{camel-oasis}、\texttt{camel-ai}、\texttt{zep-cloud}，工程命名保留了 \texttt{mirofish-backend}，前端部分页面仍然保留“图谱构建、环境搭建、模拟、报告、互动”的工作流语言。这说明项目并非从零构造了一个认知安全系统，而是在已有底座上进行问题重构。

这一点不仅体现在命名上，也体现在本地 README 和主线重构文档对系统流程的重写上。当前仓库已经把主流程明确表述为：\texttt{Input -> PrivacySanitizer -> PersonaStateVector -> RiskGraphBundle -> WorldState -> Fork A/B -> Counterfactual Risk -> Intervention}。这个表达非常关键，因为它意味着项目已经把原先更宏观的 MiroFish 方法压缩成若干可执行对象：输入不再直接进入“世界模拟”，而是先经过脱敏、画像和风险图谱构造，再由 \texttt{MiroFishRuntime} 初始化 \texttt{WorldState} 并在关键不可逆节点触发 \texttt{Fork}。

更具体地说，当前项目从 MiroFish 继承的并不是“群体规模”本身，而是三种更适合保留到竞赛版中的能力。第一是\textbf{状态化建模能力}，也就是把原始文本、上下文和角色关系压缩为可运行对象；第二是\textbf{分叉式推演能力}，也就是在关键节点上推进两条不同路径并比较结果；第三是\textbf{报告化输出能力}，也就是将复杂中间过程组织成前端和文档可以消费的结构化结论。在当前项目里，这三种能力分别落在 \texttt{PersonaStateVector / RiskGraphBundle}、\texttt{WorldState + ForkComparison} 和 \texttt{CounterfactualReporter} 上。

更关键的是，当前服务编排中已经出现 \texttt{MiroFishRuntime}、\texttt{world\_state\_snapshot}、\texttt{fork\_comparison} 等对象。这意味着，MiroFish 的状态演化和分叉思路已经不是外部叙事，而是实际进入了 CogSec 主链之中。

\subsection{为什么个人决策图谱比直接复用原始叙事更适合竞赛}

如果完全沿用原版 MiroFish 的“群体智能预测引擎”叙事，项目在信安竞赛语境中容易出现两个问题。第一，评委会把注意力放到“系统是不是一个太大的通用预测平台”，从而稀释项目在诈骗防御上的核心价值。第二，宏大的世界仿真叙事会天然抬高可验证门槛，容易让展示和代码现状之间出现不必要的张力。

相反，把项目明确收束为\CorePosition{}，则有三个优势：
\begin{enumerate}[leftmargin=2em]
  \item 它保留了 MiroFish 的底座优势，但把问题定义压缩到一个清晰、强相关的竞赛主题上；
  \item 它让系统输出从“宏大预测”变成“关键节点上的分支比较”，更容易展示也更容易答辩；
  \item 它让工程现状、页面设计、理论基础和创新点更容易互相对齐。
\end{enumerate}

\subsection{MiroFish 与当前项目的关键差异}

\begin{longtable}{p{0.20\textwidth}p{0.34\textwidth}p{0.36\textwidth}}
\toprule
\textbf{维度} & \textbf{原版 MiroFish} & \textbf{当前项目} \\
\midrule
\endfirsthead
\toprule
\textbf{维度} & \textbf{原版 MiroFish} & \textbf{当前项目} \\
\midrule
\endhead
核心对象 & 群体反应、环境变化和未来场景排演。 & 个体在诈骗场景中的关键决策路径。 \\
输入类型 & 报告、新闻、政策、故事、金融材料等现实种子。 & 对话、文本、截图、问卷和上下文描述。 \\
输出形态 & 预测报告与交互式数字世界。 & 个人决策图谱、双分支对照、可逆性窗口与干预建议。 \\
方法重心 & 图谱构建、环境搭建、大规模 Agent 模拟。 & T0/T1、风险因子向量、ThreatKnowledgeRAG、Fork、处方报告。 \\
解释单位 & 世界状态与群体演化。 & 关键动作节点、触发因子、分支差异和干预路径。 \\
\bottomrule
\end{longtable}

\section{同类竞品与 GitHub 开源调研}

\subsection{原版 MiroFish 与 MiroFish-Offline}

对于当前项目来说，最重要的开源参照首先不是反诈工具，而是原版 MiroFish 自身。根据其官方 README，MiroFish 将自己定义为一个面向现实种子材料的多智能体预测引擎：用户上传报告、新闻、政策、故事或金融材料，系统随后执行图谱构建、环境搭建、仿真执行、报告生成与深度交互\SrcMiroFishRepo。 如果把这套 README 逻辑压缩成一句话，那么原版 MiroFish 做的是“把复杂现实材料转成一个可运行的平行世界，再观察这个世界会如何演化”。

MiroFish-Offline 则保留了这一套方法逻辑，但把它改造成更适合本地运行的离线分支。它延续图谱、环境、模拟和报告的主结构，只是在基础设施上更强调本地 Neo4j、Ollama 和离线界面。这一点对于当前项目很有启发，因为它说明 MiroFish 的真正核心不是某个特定云服务，而是“图谱化输入、状态化环境、演化式推演、报告化输出”这一整套问题求解框架。

如果按项目逻辑而不是宣传口径来拆，原版 MiroFish 大体可以理解为四段。第一段是把现实种子材料转成结构化图谱，明确人物、机构、事件、关系和约束；第二段是据此初始化一个可运行环境，让不同智能体带着角色、记忆与目标进入同一世界；第三段是在世界中推进交互与状态演化，并在重要节点持续更新环境；第四段是把推演过程压缩成报告和可交互视图。这种逻辑对当前项目特别有用，因为诈骗防御里真正需要保留的恰好也是这四段，只是“世界”被缩成了“用户当前所处的局部社工情境”，“群体演化”被缩成了“个体在关键动作前的多分支演化”。

从项目逻辑上看，当前 CogSec 与原版 MiroFish 的关系，恰恰可以通过这一框架来理解：原版是把世界当作分析对象，而当前项目是把个体决策链当作分析对象。也正因为如此，在竞赛文档中把原版 MiroFish 放在“底座方法来源”而不是“完全同类竞品”的位置，反而更有利于说明项目的独立性。

\subsection{面向消费者的 AI 反诈助手}

Bitdefender Scamio 和 Norton Genie 等产品已经能够对文本、链接、网站、图片甚至深度伪造内容进行快速可疑性判断\SrcConsumerTools。 这类产品的优势在于交互简单、反馈迅速、面向普通用户部署成本低。它们解决的是“这段内容是否值得怀疑”“当前是否应该提高警惕”的问题。

然而，这类系统通常将内容本身作为分析单位，其主要输出仍是可疑性判断与通用建议，而不是围绕用户即将执行的动作构造路径比较。也就是说，它们更像是智能型的“第二意见”，而不是“围绕下一步动作的反事实决策图谱”。

\subsection{企业安全意识训练与人因风险平台}

SoSafe、KnowBe4 AIDA 与 Adaptive Security 已将个性化钓鱼演练、动态难度调整、多渠道社会工程模拟、人因风险评分和培训闭环做成较成熟的产品\SrcEnterprisePlatforms。 这一类平台已经充分证明，“个体差异”“动态风险”“个性化干预”并非空洞概念，而是现实可落地的企业安全能力。

但这些平台的核心目标是\textbf{组织层面的长期风险下降}，而不是\textbf{单个用户当前节点的即时决策辅助}。它们更关注员工未来不要再犯相同错误，而不是当前某个高危动作能否被及时改写。因此，这些产品构成了重要参照，但并不直接覆盖本项目的核心问题。

\subsection{组织级安全运营 Agent}

Microsoft Defender 的 Phishing Triage Agent 主要服务于 SOC 团队，用于帮助分析师理解用户上报的疑似钓鱼邮件并完成 triage\SrcPhishingTriage。 这表明解释性、图形化 reasoning 和自动化分析已经进入现实安全产品需求。

但此类系统的对象是组织级告警和分析师工作流，而不是普通用户的即时决策过程。它说明“解释过程可视化”具有现实价值，却并未把个体动作前的分支比较作为主要交互对象。

\subsection{防御训练型开源项目}

Pretexta 是一个很有代表性的开源社会工程训练实验室。其 README 的缩减简介可以概括为：它不是给攻击者写脚本的工具，而是一个“Train defenders, not attackers”的受控防御训练环境；系统通过 AI roleplay、自适应难度、赛后复盘和基于说服原则的分析逻辑，把社会工程学训练组织成一条完整流程\SrcPretexta。 换句话说，它要解决的问题不是“这次输入到底是不是诈骗”，而是“让使用者在模拟中看懂社会工程如何一步步推进”。

从项目逻辑上看，Pretexta 更像“训练场”而不是“决策助手”。其核心机制通常围绕三步展开：首先由 AI 角色发起对话或诱导；随后系统根据互动过程判断是否触发了危险边界、是否命中了某些说服原则；最后在复盘阶段给出赛后拆解和行为反馈。这样的逻辑对本项目有两层参考价值：一是它证明了说服原则、角色扮演和复盘输出确实可以落成工具；二是它也提醒我们，若当前项目只做到“和 AI 攻击者对话并事后点评”，那么与 Pretexta 的差异会明显缩小。

进一步说，Pretexta 代表的是“训练闭环”范式，而当前项目要争取的是“判别闭环”范式。前者更关心训练样本如何生成、参与者如何在互动中暴露弱点、复盘如何帮助其下次改进；后者更关心当下输入是否已经进入危险轨道、主 Fork 节点在哪里、哪一种保护动作还来得及把路径拉回来。两者都研究社会工程机制，但时间尺度、使用对象和核心输出并不相同。把这一点在报告中讲清楚，能够避免评委把当前项目误解为“另一个模拟训练器”。

不过，Pretexta 的主场景仍是训练与复盘，而不是基于真实输入情境生成一张实时的个人决策图谱。因此，它是很强的相邻方案，却不是与本项目完全同位的替代物。

\subsection{AXIS：研究型反事实解释工作}

AXIS 更适合被理解为研究型参照，而不是现成竞品。其核心思想，是把反事实模拟与语言模型结合起来，用于解释复杂多智能体系统为什么会产生某种行为差异\SrcAxis。 如果压缩成 README 级的理解，它回答的是：“当一个复杂系统已经做出某种行为后，怎样通过替换路径或重构条件，生成更有说服力的解释。”

从逻辑上看，AXIS 关注的是解释环节：先有行为轨迹，再做 interrogative simulation 和 counterfactual explanation，最后比较不同条件下系统输出的变化。对于当前项目，这一思路的价值不在于直接替代产品，而在于它证明了双分支比较不只是展示形式，而是一种有研究根基的解释策略。也就是说，当前项目里 Branch A / Branch B 的设计，并不是为了把报告做得更“好看”，而是为了让系统能够从“结果判断”转向“路径差异解释”。

\subsection{综合判断}

综合公开产品页、GitHub 仓库与研究工作来看，目前已有方案在内容检测、企业训练、SOC triage、训练实验室和解释研究方面都相当成熟；但围绕“实时个人反诈骗决策图谱 + 关键动作前的双分支反事实推演 + 个体风险因子命中链 + 可逆性窗口”这一组合能力，仍然存在清晰的差异化空间。也正因为如此，本项目应避免把创新点泛化成“也用 AI、也做人因风险、也能做解释”，而应更集中地强调“个人关键动作前的路径比较与干预”。

\section{为什么系统要这样设计：理论与方法依据}

\subsection{问题本质是“决策前干预”而不是“内容后判断”}

社会工程攻击之所以长期有效，不仅因为攻击文本表面上更像真的，更因为攻击者能够在合适的时机触发目标的错误行为。相关综述指出，生成式 AI 放大的并不只是内容质量，而是对人类注意、信任和判断节奏的操控能力\SrcDigitalDeception。 因此，项目将分析单位从文本转向动作，是由问题本质决定的，而不是任意选择。

\subsection{双系统理论支持 T0 / T1 分层}

双系统理论及其在电信诈骗研究中的应用表明，在高压力、高不确定性和强情绪刺激条件下，人更容易依赖快速、自动、直觉的 System 1，而较少动用缓慢、审慎的 System 2\SrcDualSystem。 这直接支持了系统的 T0 / T1 分层设计：T0 负责在毫秒级对高危动作亮红灯，为用户争取“停下来”的机会；T1 负责在更长时间窗口内恢复分析、解释与比较过程。

\subsection{用户情境与个体差异支持风险向量建模}

NIST 的用户情境研究指出，钓鱼判断的难度与用户所处的上下文强相关；而关于社会工程认知机制的研究也显示，攻击效果并不只来自内容本身，而是与个体的短期状态、长期倾向和记忆机制共同相关\SrcUserContext。 因此，系统引入多维个体风险因子，并将其表示为 \texttt{persona\_state\_vector}，并不是为了做人格标签化，而是为了描述当前场景下“为什么某个分支对这个人更危险”。

\subsection{社会学与认知心理学背景如何进入项目主线}

如果把项目仅仅理解为“识别诈骗话术”，那么很多真正决定受骗与否的因素就会被遗漏。社会工程攻击在现实中常常不是靠单一文本骗过用户，而是通过\textbf{关系结构、权威位置、信息不对称、情绪压迫和社会隔离}逐步推动行为。这也是为什么当前项目不能只停留在文本分类层，而要把社会学与认知心理学的解释变量一起纳入主链。

从社会机制角度看，攻击者通常会先构造一个不平等的互动结构：例如伪装成客服、警方、领导或熟人，先建立权威或关系优势；再通过“必须保密”“不要告诉别人”“现在马上处理”的方式切断受害者与外部核验链路；最后把用户推入一个既缺少信息、又缺少社会支持、还伴随强时限压力的局部环境。这个过程本质上不是简单的信息传递，而是\textbf{对决策环境的重构}。当前项目中 \texttt{authority\_intensity}、\texttt{info\_asymmetry}、\texttt{social\_isolation} Fork 节点以及 \texttt{help\_seeking}、\texttt{verification\_habit} 等保护性维度，正是为了把这种社会结构变化落成可计算对象。

从认知心理学角度看，真实受骗并不只是“知识不够”，而是“在高负荷、高情绪、高紧迫情境下，System 1 抢先接管了决策”。也就是说，用户明明知道应该核验，但在时间压力、损失厌恶、情绪波动和权威压迫叠加时，仍可能直接执行高风险动作。这也是为什么项目的 18 维画像没有只围绕“是否懂安全”展开，而是同时建模了\textbf{当前状态、脆弱点和保护机制}。其中状态维度回答“现在是不是正处于容易失手的时候”，脆弱性维度回答“哪类诱导最容易命中”，保护性维度回答“还有哪些认知和行为护栏能把用户拉回来”。

从行为机制角度看，当前项目最想解决的是“风险如何一步步推进”。例如，当用户先被权威压制，再被要求保密，随后被诱导共享屏幕或提供验证码时，系统需要说明：风险不是在最后一步才突然出现，而是在前面几个节点就已经开始累积；而一旦 \texttt{help\_seeking} 被切断、\texttt{verification\_habit} 被压制、\texttt{time\_pressure} 与 \texttt{emotional\_volatility} 同时升高，系统就会更倾向于判定其进入 \texttt{SYSTEM\_1} 主导状态。也正因为如此，当前项目并没有把社会学背景和心理学背景放在文献综述页就结束，而是把它们写进了 \texttt{PersonaStateVector}、\texttt{RiskGraphBundle} 和 \texttt{ForkComparison} 这些运行时对象之中。

\subsection{本地实现中的 18 维特征分组与运行时映射}

这一点在本地项目中不是抽象设想，而是已经被明确写成了可执行结构。\texttt{cognitive\_profiler.py} 中的 18 维画像被直接分成三类。第一类是\textbf{状态特征}：\texttt{time\_pressure}、\texttt{financial\_pressure}、\texttt{info\_asymmetry}、\texttt{emotional\_volatility}、\texttt{cognitive\_load} 和 \texttt{authority\_intensity}，对应用户当下处在什么样的压力与情境中。第二类是\textbf{脆弱性特征}：\texttt{authority\_compliance}、\texttt{social\_proof\_sensitivity}、\texttt{scarcity\_sensitivity}、\texttt{loss\_aversion\_threshold}、\texttt{gambler\_fallacy} 和 \texttt{trust\_threshold}，对应哪些攻击机制更容易被命中。第三类是\textbf{保护性特征}：\texttt{decision\_delay}、\texttt{verification\_habit}、\texttt{help\_seeking}、\texttt{link\_check\_ability}、\texttt{transaction\_review} 和 \texttt{prior\_experience}，对应用户在危险时是否具备延迟、核验和求助能力。

这一分组与当前项目的理论主张是对齐的。状态特征决定的是此刻是否已进入高压、高负荷、强不对称情境；脆弱性特征决定的是哪类说服原则和攻击策略更可能生效；保护性特征决定的是系统还有没有机会把路径从 Branch A 拉回 Branch B。也就是说，项目并没有把所有心理变量混成一个总分，而是明确区分了“此刻的状态”“容易被什么击穿”“还有什么能把人拉回来”三层逻辑。

更进一步，\texttt{CognitiveProfile.to\_persona\_state\_vector()} 已经把这 18 维画像压缩为运行时数字孪生向量。比如，\texttt{analytic\_control} 来自 \texttt{decision\_delay + link\_check\_ability + transaction\_review} 的均值，\texttt{fomo\_susceptibility} 来自 \texttt{scarcity\_sensitivity + social\_proof\_sensitivity} 的组合，\texttt{digital\_trust\_boundary} 用信任阈值和链接核验能力共同计算，\texttt{risk\_recovery\_awareness} 则由既往经验、求助倾向和核验习惯共同构成。当前代码还显式把 \texttt{time\_pressure > 7 且 emotional\_volatility > 7} 作为 \texttt{SYSTEM\_1} 触发条件，并在 \texttt{switch\_policy} 中保留了“核验习惯与决策延迟较高时允许回到 System 2”的重入规则。这使“画像”真正落成了“运行时状态”，而不只是文案标签。

与此同时，当前抽取器也并非只输出分数。它还在本地代码中显式维护了十类场景关键词表，包括虚假征信、刷单返利、虚假网络投资理财、冒充电商物流客服、冒充公检法及政府机关、冒充领导熟人、网络游戏产品虚假交易、网络婚恋交友、机票退改签以及虚假贷款代办信用卡等常见类别；并通过轻量 \texttt{LIWC\_DIMENSION\_MAP} 将焦虑、恐惧、紧迫、金钱、权威等语言信号映射到画像维度上，用作交叉校验。这使当前实现既保留了社会工程学与语言心理学的解释视角，又没有把系统写成一个不可控的纯文本情绪分析器。

从理论对应关系上看，这 18 维并不是随意拼接出来的。状态特征更接近双系统理论里“当前是否已进入高压、低审慎加工”的触发条件；脆弱性特征更接近说服原则研究中“哪类刺激更容易命中个体弱点”的差异项；保护性特征则更接近用户情境研究和行为安全研究中“是否存在延迟、核验、求助、复盘经验”等保护机制。也正因为采用了这种三层分组，当前项目才有能力把“为什么危险”解释成三条并行链：此刻的状态为什么糟、攻击策略为什么容易命中、系统还有哪些抓手能把人拉回安全分支。

\subsection{18 维画像维度总表}

\begin{longtable}{p{0.18\textwidth}p{0.20\textwidth}p{0.26\textwidth}p{0.24\textwidth}}
\toprule
\textbf{类别} & \textbf{维度} & \textbf{在项目中的含义} & \textbf{主要对应的理论背景} \\
\midrule
\endfirsthead
\toprule
\textbf{类别} & \textbf{维度} & \textbf{在项目中的含义} & \textbf{主要对应的理论背景} \\
\midrule
\endhead
状态特征 & \texttt{time\_pressure} & 用户是否被压入“必须马上处理”的决策节奏。 & 双系统理论中的快速加工触发；紧迫/稀缺说服原则。 \\
状态特征 & \texttt{financial\_pressure} & 经济焦虑是否放大了对退款、挽损、赚钱承诺的敏感性。 & 损失厌恶、压力下风险偏移、诈骗中的收益/损失 framing。 \\
状态特征 & \texttt{info\_asymmetry} & 用户是否处于对平台规则、资金状态、事件真相不透明的位置。 & 用户情境研究；社会工程中的信息控制与结构性不对称。 \\
状态特征 & \texttt{emotional\_volatility} & 情绪波动是否已干扰冷静核验。 & 恐惧/焦虑启动、情绪对启发式判断的放大作用。 \\
状态特征 & \texttt{cognitive\_load} & 用户当前是否处于多任务、信息过载或理解负担过高状态。 & 认知负荷理论；高负荷下 System 2 介入不足。 \\
状态特征 & \texttt{authority\_intensity} & 当前情境中的权威压迫感有多强。 & 权威服从、社会角色与权力不对称。 \\
脆弱性特征 & \texttt{authority\_compliance} & 用户对“官方/警方/平台/领导”话语的服从倾向。 & 权威原则；社会角色服从机制。 \\
脆弱性特征 & \texttt{social\_proof\_sensitivity} & 用户是否容易被“别人都这样做”“大家都在用”影响。 & 社会认同/从众效应。 \\
脆弱性特征 & \texttt{scarcity\_sensitivity} & 对名额、时限、窗口关闭等稀缺信号是否敏感。 & 稀缺原则；FOMO 心理。 \\
脆弱性特征 & \texttt{loss\_aversion\_threshold} & 面对损失威胁时是否会优先保住眼前利益而跳过核验。 & 损失厌恶；风险规避与挽损心态。 \\
脆弱性特征 & \texttt{gambler\_fallacy} & 是否容易相信“再操作一次就能追回”“再跟一步就能翻盘”。 & 赌徒谬误；升级承诺。 \\
脆弱性特征 & \texttt{trust\_threshold} & 对陌生联系人、链接和流程的信任边界是否偏低。 & 数字信任边界；关系伪装与熟悉性偏差。 \\
保护性特征 & \texttt{decision\_delay} & 用户是否具备“先停一下再决定”的能力。 & 延迟决策、冲动控制、System 2 重入。 \\
保护性特征 & \texttt{verification\_habit} & 遇到异常信息时是否会主动二次核验。 & 安全行为习惯；NIST 用户情境中的防御行为。 \\
保护性特征 & \texttt{help\_seeking} & 用户是否倾向于向家人、同事或官方渠道求助。 & 社会支持网络；反隔离机制。 \\
保护性特征 & \texttt{link\_check\_ability} & 是否具备识别可疑链接、域名和页面异常的能力。 & 基础安全素养；线索甄别能力。 \\
保护性特征 & \texttt{transaction\_review} & 转账或授权前是否会复核对象、金额与流程。 & 关键动作前复核；交易安全行为。 \\
保护性特征 & \texttt{prior\_experience} & 过去经历是否帮助用户更早识别套路。 & 经验学习；风险恢复意识与情境记忆。 \\
\bottomrule
\end{longtable}

这张表之所以重要，是因为它把“理论背景”和“项目实现”对齐了。报告里提到的社会支持、权威压力、情绪波动、信息不对称、核验习惯和求助机制，并不是散落在不同论文里的概念，而是已经分别落成了项目中的显式维度。也正因为有了这套维度表，后续的 \texttt{PersonaStateVector} 压缩、\texttt{persona\_weakness\_hits} 命中链、\texttt{Fork} 选择与 \texttt{InterventionPrescription} 生成才会显得顺理成章，而不是一串彼此断开的模块。

\subsection{说服原则支持攻击-人-环境图谱}

无论是学术研究，还是 Pretexta 这类开源训练系统，都持续强调社会工程攻击中说服原则的核心作用\SrcPersuasion。 权威、紧迫、承诺一致性、社会认同、互惠和喜好，不只是攻击话术的语言表面，更是推动用户跨过关键动作节点的行为触发机制。正因如此，系统将这些因素映射进攻击-人-环境图谱，并用于解释某一分支为何在当前时刻快速增险。

\subsection{反事实解释支持 Branch A / Branch B 设计}

AXIS 等研究说明，复杂行为系统的解释不仅可以通过“发生了什么”来实现，也可以通过“如果换一条路径会怎样”来实现\SrcAxis。 对本项目而言，这一思路非常关键。因为用户并不总是需要知道“系统为何判断危险”，他们更需要知道“我如果继续做会怎样”“我如果此刻停下来会怎样”。因此，Branch A / Branch B 不是展示形式上的锦上添花，而是系统解释性的核心载体。

\subsection{本地代码中的主链时序与对象关系}

如果把项目逻辑完全按本地实现展开，当前主链并不是若干独立模块的松散拼接，而是一条由 \texttt{CogSecService.analyze\_text()} 统一编排的运行时时序。其一，输入文本先进入 \texttt{PrivacySanitizer.sanitize\_with\_retry()}，并以 \texttt{session\_only} 作为默认存储策略；其二，\texttt{T0FastResponder.scan()} 给出亚秒级红旗护栏；其三，\texttt{CognitiveProfileExtractor.extract()} 输出 18 维画像，再由 \texttt{to\_persona\_state\_vector()} 生成数字孪生；其四，\texttt{ThreatKnowledgeRAG.build\_risk\_graph\_bundle()} 生成 \texttt{risk\_graph\_bundle}；其五，\texttt{MiroFishRuntime.run()} 先初始化 \texttt{WorldStateSnapshot(step=0)}，再选择主 \texttt{Fork} 节点并同时推进 Branch A / Branch B；其六，\texttt{RiskScorer.evaluate\_counterfactual()} 不是对原始文本打分，而是对两条分支的 \texttt{trajectory\_gap}、\texttt{irreversibility\_loss}、\texttt{evidence\_consistency} 和最终风险进行聚合；其七，\texttt{CounterfactualReporter.generate\_report()} 将这些对象重组为数字孪生摘要、弱点链、攻击链、分支对照、不可逆节点和干预处方。

这一时序说明了一个特别重要的事实：当前项目的“最终风险”并不是由某个单独分类器对输入文本打出来的，而是由 \texttt{persona\_state\_vector}、\texttt{risk\_graph\_bundle}、\texttt{world\_state\_snapshot}、\texttt{branch\_a\_log / branch\_b\_log} 和 \texttt{fork\_comparison} 共同决定的。换句话说，它确实已经在工程上形成了“先理解人和环境，再构造路径，再比较后果”的主判别逻辑。

\subsection{为什么采用 T0 + 风险向量 + RAG + Fork + Reporter 的组合}

结合以上理论与问题约束，系统当前采用的模块组合具有明确原因：
\begin{enumerate}[leftmargin=2em]
  \item \textbf{T0 快速护栏} 对应高危红旗的即时阻断需求；
  \item \textbf{个体风险向量} 对应个体差异和用户情境的解释需求；
  \item \textbf{ThreatKnowledgeRAG} 对应知识约束和证据锚定需求；
  \item \textbf{Fork} 对应关键动作节点上的路径分叉需求；
  \item \textbf{CounterfactualReporter} 对应可视化、可解释和可执行建议输出需求。
\end{enumerate}

换句话说，这一系统设计并不是功能堆叠，而是从“项目起点”“竞品空位”和“理论支持”共同推导出来的结果。

\section{项目定位与总体目标}

\subsection{项目定位}

\ProjectName{} 的核心定位是\CorePosition{}。它强调三点：
\begin{enumerate}[leftmargin=2em]
  \item 分析对象不是单条消息，而是围绕关键动作展开的个人决策路径；
  \item 系统输出不是单一风险分，而是分支差异、可逆性窗口和个性化干预建议；
  \item 工程形态不是纯研究框架，而是具备前端工作台、后端编排和报告输出的完整系统。
\end{enumerate}

换句话说，项目既可以被理解为一个认知安全分析系统，也可以被理解为一个“面向高风险动作前一刻”的个人辅助判断工具。它所服务的不是事后复盘，而是事前干预；它所追求的也不是抽象预测，而是帮助用户在关键节点前做出更安全的选择。

\subsection{目标问题}

项目试图回答的核心问题是：
\begin{quote}
当用户在社会工程学攻击场景中面临关键动作决策时，系统如何在损失真正发生之前，识别高风险节点、解释其形成机制，并通过反事实推演给出更安全的替代路径。
\end{quote}

\subsection{应用场景}

系统当前重点面向以下几类场景：
\begin{enumerate}[leftmargin=2em]
  \item 冒充平台客服、物流退款、订单异常、售后理赔等高频诈骗场景；
  \item 冒充公检法、上级、熟人或合作方实施的高压诱导场景；
  \item 索取验证码、屏幕共享、远程控制、转入安全账户等高危动作场景；
  \item 企业员工面对钓鱼邮件、身份伪装、权限请求时的决策辅助场景；
  \item 适用于展示、培训与答辩演示的结构化可视化场景。
\end{enumerate}

\section{系统总体架构}

\subsection{总体设计思路}

系统采用“输入治理层—分析理解层—反事实推演层—报告与干预层”的四层结构：
\begin{enumerate}[leftmargin=2em]
  \item \textbf{输入治理层}：负责接收文本、聊天记录、截图 OCR、问卷信息与上下文描述，并进行隐私脱敏与格式规整；
  \item \textbf{分析理解层}：负责个体风险因子提取、场景识别、威胁知识检索、攻击-人-环境图谱构建；
  \item \textbf{反事实推演层}：负责在关键动作节点执行 Fork，对继续执行与暂停核验两种分支进行比较；
  \item \textbf{报告与干预层}：负责输出数字孪生概览、分支对照、可逆性曲线、干预窗口与个性化处方。
\end{enumerate}

\subsection{主链工作流}

系统主链可以概括为以下步骤：
\begin{enumerate}[leftmargin=2em]
  \item 输入场景文本、上下文和可选问卷；
  \item 执行隐私脱敏和 T0 快速护栏扫描；
  \item 提取多维个体风险因子，生成 \texttt{persona\_state\_vector}；
  \item 通过 \texttt{ThreatKnowledgeRAG} 构造攻击-人-环境风险图谱；
  \item 由 \texttt{MiroFishRuntime} 初始化 \texttt{world\_state\_snapshot}；
  \item 在高危不可逆节点执行 Fork，生成 Branch A / Branch B 轨迹；
  \item 使用 \texttt{RiskScorer} 聚合分支差值；
  \item 由 \texttt{CounterfactualReporter} 输出结构化报告和干预处方。
\end{enumerate}

\subsection{主链对象关系}

从对象关系而不是页面顺序来理解，当前系统实际上围绕五层核心对象运转。第一层是 \texttt{CognitiveProfile}，它承载 18 维画像及其置信度；第二层是 \texttt{PersonaStateVector}，它把画像压缩成可运行的数字孪生；第三层是 \texttt{RiskGraphBundle}，它将攻击策略、说服原则、证据项、弱点命中链、资产目标和 Fork 候选组织成一张可计算风险图谱；第四层是 \texttt{WorldStateSnapshot + BranchTraceStep + ForkComparisonResult}，它们共同承载运行时状态、分支轨迹和主分叉比较；第五层是 \texttt{CFReport / InterventionPrescription}，它们负责把运行时对象转译成用户能读懂、前端能展示的解释性报告。

这组对象关系非常重要，因为它说明当前项目的结构不是“检测器 + 报告模板”，而是“画像对象 + 图谱对象 + 运行时对象 + 报告对象”的层层收束。也正因为如此，前端五屏并不是五个彼此独立的页面，而是同一份主链结果在不同视角下的投影。

\subsection{T0 / T1 双层机制}

系统显式采用 T0 / T1 双层架构：
\begin{enumerate}[leftmargin=2em]
  \item \textbf{T0 快速响应层}：面向验证码、转账、安全账户、远程控制、屏幕共享等强危险指令做毫秒级检测，承担红线护栏职责；
  \item \textbf{T1 深度推演层}：面向复杂场景做画像、图谱、Fork 比较和反事实报告生成，承担解释与干预职责。
\end{enumerate}

\section{核心技术路线}

\subsection{个体风险因子与数字孪生}

系统通过 \texttt{CognitiveProfileExtractor} 提取与行为决策相关的多维风险因子，并将其压缩为 \texttt{persona\_state\_vector}。这一向量不是为了做永久人格判断，而是为了在当前场景下刻画决策脆弱性结构。它使系统能够回答“哪一类触发机制最有效”“此人更易在哪种节点失守”“哪些保护性习惯能够将路径拉回安全分支”。

从本地实现上看，这一层并不是简单把 18 个分数原样透传到前端，而是做了两层压缩。第一层是 \texttt{overall\_vulnerability\_score()} 对三类特征进行加权聚合：脆弱性特征权重更高、状态特征次之、保护性特征用于削减风险。第二层是 \texttt{to\_persona\_state\_vector()} 将若干相关维度合成为运行时变量，例如 \texttt{analytic\_control}、\texttt{fomo\_susceptibility}、\texttt{digital\_trust\_boundary}、\texttt{asset\_sensitivity}、\texttt{risk\_recovery\_awareness}、\texttt{system1\_bias} 和 \texttt{system2\_control}。这说明项目并没有把“画像”理解成静态画像卡，而是把它视为驱动后续 WorldState 和 Fork 演化的状态初始化层。

\subsection{ThreatKnowledgeRAG 与风险图谱}

\texttt{ThreatKnowledgeRAG} 负责在威胁知识、案例证据和当前输入之间建立约束关系。与纯生成式系统相比，这一模块的意义在于让推演尽量锚定已知诈骗模式、诱导链和红旗动作，降低悬浮生成和无依据解释的风险。输出的 \texttt{risk\_graph\_bundle} 不只是可视化图谱，更是后续 Fork 和报告生成的证据基础。

从代码逻辑看，\texttt{ThreatKnowledgeRAG} 内部首先以 \texttt{FraudCase} 和 \texttt{AttackStrategy} 作为知识单元，围绕场景类型和画像状态检索最相关的攻击策略；随后再推导 \texttt{persona\_weakness\_hits}、\texttt{evidence\_items}、\texttt{asset\_targets}、\texttt{environment\_context} 和 \texttt{fork\_points}，最终组装成 \texttt{RiskGraphBundle}。其中 \texttt{persuasion\_principles} 会直接收集进入图谱，\texttt{consistency\_score} 会用于衡量当前图谱是否足够可信；一旦低于阈值，系统还会触发 \texttt{hallucination\_rollback} 并回退到已验证的威胁模板。也就是说，这一层不是“检索几个案例看看”，而是主链中承担知识约束、证据组织和 Fork 候选生成的核心模块。

\subsection{Fork 与反事实双分支}

系统将关键动作节点视为 Fork 点，并围绕该节点同时构造两条路径：
\begin{enumerate}[leftmargin=2em]
  \item \textbf{Branch A}：继续执行当前诱导动作，如继续转账、继续共享屏幕、继续提供验证码；
  \item \textbf{Branch B}：暂停、核验、延时、回拨、咨询可信第三方或直接中止操作。
\end{enumerate}

双分支设计使风险分析从静态评分转变为路径比较。系统真正关注的不是某一时刻的绝对风险，而是\textbf{两条路径在未来几步上的差异、转折点与可逆性变化}。

当前本地实现已经把主 Fork 类型具体写成六类：\texttt{transfer\_money}、\texttt{screen\_share}、\texttt{verification\_code}、\texttt{unknown\_app\_download}、\texttt{social\_isolation} 和 \texttt{fake\_official\_verification}。这几个类型很能说明项目思路：它不仅关注资金和凭证泄露，也关注“不要告诉别人”“单独联系”这类切断社会支持系统的社工步骤。与此同时，\texttt{initialize\_world\_state()} 已经会为每次输入生成 \texttt{trust\_score}、\texttt{asset\_exposure}、\texttt{intervention\_window}、\texttt{posterior\_risk} 和 \texttt{reversibility} 等初始状态字段，说明 Fork 前后并不是抽象故事，而是被压进了统一的运行时状态空间。

\subsection{风险聚合与报告输出}

\texttt{RiskScorer} 根据分支轨迹、后验更新、可逆性曲线、证据一致性与异常标记聚合整体风险；\texttt{CounterfactualReporter} 则将这些中间对象转译成用户可理解的结构化结果，包括：
\begin{enumerate}[leftmargin=2em]
  \item 当前场景与个体状态摘要；
  \item 关键动作节点及其危险原因；
  \item Branch A / Branch B 的分步对照；
  \item 最佳干预窗口；
  \item 个性化处方与建议动作。
\end{enumerate}

更细一点看，\texttt{RiskScorer.evaluate\_counterfactual()} 实际聚合的不是一条单一风险曲线，而是 \texttt{persona\_prior}、\texttt{attack\_match}、\texttt{trajectory\_gap}、\texttt{irreversibility\_loss}、\texttt{evidence\_consistency} 五类量，并进一步推导 \texttt{low\_discriminability}、\texttt{audit\_required} 等状态。对应地，\texttt{CounterfactualReporter.generate\_report()} 也不仅输出摘要文字，而会把数字孪生摘要、认知弱点链、攻击策略链、Fork 节点、Branch 对照、不可逆节点、最佳干预窗口和 \texttt{intervention\_prescriptions} 一并组织进结构化报告。这些实现细节说明，项目的结果层确实已经开始从“解释文本”转向“解释对象化”。

\subsection{可解释性的工程意义}

对本项目而言，可解释性不是附属功能，而是主输出之一。用户、评委或分析者不应只看到“风险 82 分”，而应能看到“风险为何在第 2 至第 3 步快速升高”“哪些触发因子导致进入 Branch A”“在哪个节点仍有可逆空间”“哪一条建议与当前脆弱性结构最匹配”。

\section{工程实现现状}

\subsection{后端实现现状}

当前后端采用 Python 3.11+ 和 Flask 作为服务基础，依赖中已包含 \texttt{openai}、\texttt{camel-oasis}、\texttt{camel-ai}、\texttt{chromadb}、\texttt{PyMuPDF}、\texttt{presidio-analyzer} 等模块，说明系统在 LLM 调用、多智能体生态、知识检索、文件处理和隐私治理方面都已具备扩展基础。

\texttt{run.py} 体现了较标准的应用工厂启动结构；\texttt{run\_cogsec\_demo.py} 提供了独立的 CogSec 最小 Demo 服务，暴露页面、健康检查、样例输入和分析接口。这种“双入口”设计具有明确工程意义：一方面保留完整系统扩展空间，另一方面保证最小演示链路的稳定性。

如果继续往细处看，后端里的两个“前置模块”已经写得比较扎实。其一，\texttt{T0FastResponder} 不是泛泛的“快筛模块”，而是一个目标亚秒级返回的正则护栏；当前默认就内置了 \texttt{screen\_share\_bank\_card}、\texttt{safe\_account\_transfer} 和 \texttt{police\_secrecy\_isolation} 三类高危规则，分别覆盖共享屏幕+银行卡、所谓安全账户转账，以及“警方/公安 + 保密/隔离联系”的典型社工组合。其二，\texttt{PrivacySanitizer} 也不是只做单次替换，而是优先调用 Presidio、失败则回退到本地正则，并对身份证号、手机号、银行卡和姓名等敏感实体做 \texttt{session\_only} + \texttt{minimal\_necessary} 脱敏；如果首次脱敏后仍检测到 PII，系统还会自动再重试一次。也就是说，后端在真正进入画像和推演前，已经具备了较明确的护栏与治理前处理层。

\subsection{主链服务编排现状}

服务编排层中已经存在 \texttt{CogSecService}，并直接组合使用 \texttt{PrivacySanitizer}、\texttt{T0FastResponder}、\texttt{CognitiveProfileExtractor}、\texttt{ThreatKnowledgeRAG}、\texttt{MiroFishRuntime}、\texttt{RiskScorer} 和 \texttt{CounterfactualReporter}。这说明系统已不只是“前端演示 + 零散脚本”，而是具备明确的后端主链结构。

更重要的是，当前结果对象已包含 \texttt{persona\_state\_vector}、\texttt{risk\_graph\_bundle}、\texttt{world\_state\_snapshot}、\texttt{fork\_comparison}、\texttt{intervention\_prescriptions} 等字段。也就是说，“个人决策图谱”并不是停留在概念层，而已经映射为具体的数据结构。

这一层还有一个经常被忽略、但对竞赛展示很重要的优点：\texttt{CogSecAnalysisResult} 已经把前端、报告和评测需要看的信息统一收口了。它除了主链对象本身，还同步返回 \texttt{t0\_fast\_response}、\texttt{sanitization}、\texttt{metrics}、\texttt{anomalies} 和 \texttt{implementation\_status}。其中 \texttt{metrics} 已经预留了 \texttt{T0 latency < 500ms}、\texttt{End-to-end inference < 30s}、\texttt{FSA > 70\%}、\texttt{CPA > 60\%}、\texttt{EWT > 2 steps}、\texttt{FPR < 15\%}、\texttt{CAL < 0.2} 和 \texttt{EAR > 65\%} 等 benchmark 钩子；\texttt{anomalies} 则会显式汇总 \texttt{pii\_leak\_detected}、\texttt{hallucination\_rollback}、\texttt{timeout\_degraded}、\texttt{audit\_required} 等异常标志。对比赛版来说，这意味着项目已经开始把“能分析”与“能审计、能验收、能分阶段交付”放进同一条主链，而不是把它们放在文档外另算。

\subsection{最小运行时与回退设计}

\texttt{cogsec\_minimal\_runtime.py} 说明系统已具备较好的最小闭环能力。即便不依赖完整主链，也能使用风险因子提取、案例回退、图谱构造、分支打分和报告生成完成一次基本分析。这一点对竞赛尤其重要，因为它说明系统具有较好的可演示性和可复现性。

这种最小运行时思路还有一个现实价值：它把“完整平台能力”和“比赛可演示能力”明确拆开了。也就是说，即使 Zep、完整多智能体上下文或更长时域搜索暂时不上线，项目仍然可以依靠最小闭环把输入、脱敏、画像、图谱、Fork 和报告跑通。对于信安竞赛而言，这种策略往往比空喊“以后还能扩展到更大的世界模拟”更有说服力，因为它直接对应“今天能稳定展示什么”。

\subsection{前端实现现状}

前端采用 Vue 3、Vue Router、Axios、D3 和 ECharts，使用 Vite 作为开发与构建工具。技术栈本身非常适合本项目的交互形态，因为它天然支持多页面工作流、结构化图谱、曲线、时间线和报告可视化。

从当前组件结构看，系统已具备完整工作台雏形。尤其是 \texttt{CogSecWorkbench.vue} 中的五屏设计，与项目定位高度一致：
\begin{enumerate}[leftmargin=2em]
  \item 用户数字孪生；
  \item 攻击-人-环境图谱；
  \item 双分支反事实推演；
  \item 风险 / 可逆性曲线；
  \item 个性化干预处方。
\end{enumerate}

\subsection{测试与基准}

测试目录中已包含针对 T0、隐私脱敏、画像提取、ThreatKnowledgeRAG、RiskScorer、CounterfactualReporter 与主链运行时的测试文件；同时，\texttt{benchmark\_cogsec.py} 也已设置 T0 延迟和端到端耗时的基准测试钩子，目标为 T0 小于 500ms、端到端小于 30s。这表明系统已具备进一步走向正式实验评估的基础。

如果按当前仓库结构展开，这一部分已经不只是“留了测试目录”。本地 \texttt{tests/} 中至少已有 \texttt{test\_t0\_fast\_responder.py}、\texttt{test\_privacy\_sanitizer.py}、\texttt{test\_cognitive\_profiler.py}、\texttt{test\_threat\_rag.py}、\texttt{test\_risk\_scorer.py}、\texttt{test\_cf\_reporter.py} 和 \texttt{test\_mainline\_runtime.py}。这些测试已经覆盖了几个很关键的主张，例如：\texttt{PersonaStateVector.storage\_policy["default"] == "session\_only"}、图谱一致性不足时 \texttt{hallucination\_rollback} 会被触发，以及当两条分支过于接近时，评分器会将结果标记为 \texttt{low\_discriminability}。与此同时，\texttt{benchmark\_cogsec.py} 输出的也不是一条简单时间统计，而是会生成包含 \texttt{t0\_avg\_latency\_ms}、\texttt{end\_to\_end\_seconds}、\texttt{peak\_memory\_mb} 及其对应阈值判断的 JSON 结果。这些细节说明当前项目已经在往“可测、可复现、可验收”的方向走，而不是停留在概念演示层。

\subsection{当前实现状态判断}

综合工程证据，当前状态可归纳如下：
\begin{enumerate}[leftmargin=2em]
  \item \StatusDone{} 前后端分离骨架和联调路径已经建立；
  \item \StatusDone{} 存在主服务入口与最小 Demo 入口；
  \item \StatusDone{} 核心主链对象与五屏工作台均已出现；
  \item \StatusDone{} 核心模块已有测试与 benchmark 钩子；
  \item \StatusRisk{} 首页品牌与外层流程叙事仍保留部分原版 MiroFish 的泛化表述；
  \item \StatusRisk{} 正式实验结果、误报漏报案例与专家评分表仍需补齐；
  \item \StatusPlan{} 更深层的大规模 Agent 世界能力可作为底座扩展，但不宜作为比赛版主卖点。
\end{enumerate}

\section{创新点与差异化价值}

\subsection{从内容识别转向决策建模}

项目最根本的创新不在于“也使用 AI”，而在于把反诈问题的最小分析单位从消息内容转为用户决策动作。系统不只问“这段话是不是诈骗”，而是问“这个人在当前情境下接下来会不会做出危险动作”。

\subsection{个体风险因子进入路径推演}

项目并不将所有用户视为同质个体，而是通过多维风险因子向量解释不同用户在相同攻击刺激下为何会产生不同决策倾向。这样，系统给出的建议就不再只是通用提醒，而能与个体脆弱性结构产生关联。

更具体地说，这里的创新不只是“有一套画像”，而是形成了一条从画像到处方的连续对象链：\texttt{18 维 CognitiveProfile -> PersonaStateVector -> persona\_weakness\_hits -> ForkComparison -> InterventionPrescription}。也就是说，项目并不是先抽一套心理标签，再在最后手工拼几条建议；它是在运行时真正用这些特征去决定哪类攻击策略更匹配、哪类 Fork 节点更危险、哪一个干预窗口还有效、哪一种处方与当前弱点链最相关。这种“节点级命中 -> 分支级比较 -> 处方级输出”的连续性，是当前方案与普通风险说明器之间一个非常关键的区别。

\subsection{反事实报告作为主输出}

项目真正主打的不是“风险分”，而是“分支差异”。通过 Branch A / Branch B 的双路径对比，系统可以展示危险如何积累、在哪个节点开始不可逆、哪一类替代动作能够改变结果。这使输出天然具备解释性、故事性和教育性。

另外，当前 Fork 设计里显式出现 \texttt{social\_isolation} 和 \texttt{fake\_official\_verification} 两类节点，也让项目的创新不止停留在“防转账、防验证码”这种金融操作层，而是开始把社工攻击中的隔离受害者、伪造官方核验、切断社会支持系统等更偏社会机制的步骤纳入主判别链。这一点其实很重要，因为真实诈骗往往先通过隔离和权威塑造夺走核验能力，再推进资金或凭证层动作；若系统只盯着最后一跳资金动作，就很难体现其“认知安全”而非“交易风控”的特征。

\subsection{T0 / T1 双层防御结构}

T0 负责高危红旗的毫秒级响应，T1 负责复杂场景的深度解释和反事实推演。这种双层设计在防御系统中具有现实意义，因为它兼顾实时阻断与深入说明，不把所有任务都压给单次大模型调用。

\subsection{创新点成立条件}

上述创新点要真正成立，必须满足以下条件：
\begin{enumerate}[leftmargin=2em]
  \item 页面和报告中必须清楚展示关键节点、双分支和可逆性，而不能只停留在口头表述；
  \item 风险因子向量应服务于决策解释，而不是被写成夸张的人格诊断；
  \item 系统需要有足够的样例、测试和 benchmark 证明链路稳定可运行；
  \item 文档与答辩口径必须统一，不再让“大规模世界模拟”盖过“个人决策图谱”。
\end{enumerate}

\section{风险治理、隐私与伦理边界}

\subsection{治理必要性}

认知安全系统天然会接触手机号、银行卡、验证码、关系链、对话记录、身份信息和情绪表达等敏感内容。若没有治理边界，系统越擅长理解个体，越可能带来更大的隐私、误用和伦理风险。因此，风险治理必须内嵌进系统架构，而不能作为附属说明。

\subsection{当前治理基础}

项目依赖中已引入 \texttt{presidio-analyzer}，服务中也存在 \texttt{PrivacySanitizer} 和 \texttt{storage\_policy="session\_only"} 等字段；与此同时，现代安全产品也越来越强调审计、透明 reasoning、最小必要处理和用途边界\SrcGovernance。 这使项目具备将治理写成正式结构的基础。

更重要的是，这些治理点在当前仓库里并不是一句口号，而是已经落在对象字段和异常标志上。脱敏结果对象里明确包含 \texttt{storage\_policy}、\texttt{minimal\_necessary}、\texttt{retry\_count} 和 \texttt{pii\_leak\_detected}；主服务会在脱敏重试后仍检测到泄露时直接抛出异常，拒绝继续进入后续分析；风险图谱层在一致性不足时会触发 \texttt{hallucination\_rollback}；评分层和服务层又会把 \texttt{low\_discriminability}、\texttt{audit\_required}、\texttt{timeout\_degraded} 等状态汇总到 \texttt{anomalies}。换句话说，当前项目已经开始把“不能确定”“需要审计”“知识一致性不足”“隐私处理失败”这些负面状态显式编码进主链，而不是默认所有输入都能被顺畅分析。

\subsection{建议遵循的边界}

建议系统在文档与实现中坚持以下边界：
\begin{enumerate}[leftmargin=2em]
  \item 仅用于防御与决策辅助，不输出进攻性社工脚本；
  \item 优先进行本地或前置脱敏，尽量减少敏感信息暴露；
  \item 个体风险因子仅用于当前场景分析，不形成惩罚性标签；
  \item 对证据不足的情况返回“需要人工核验”，而非强行确定；
  \item 关键判断应附带证据来源、分支依据和审计信息。
\end{enumerate}

\section{当前局限与迭代路线}

\subsection{当前主要局限}

尽管系统已经具备较清晰的工程骨架与主链对象，但当前仍存在三类需要正视的问题：
\begin{enumerate}[leftmargin=2em]
  \item 外层品牌和部分页面流程仍保留原版 MiroFish 的泛化叙事，主卖点需要进一步统一；
  \item 正式实验结果、误报漏报案例和专家评分表仍需补齐，才能形成更强的量化说服力；
  \item 更深层的大规模 Agent 世界能力虽可作为底座拓展，但对比赛版主问题的直接贡献仍有限。
\end{enumerate}

除此之外，还有两类更细的工程局限也值得在竞赛文本中坦诚写出。第一，当前 T0 层仍以规则命中为主，虽然速度快、护栏明确，但覆盖面仍偏向少数高危硬规则，后续若要提高泛化性，仍需要更多场景样本和规则/模型协同。第二，当前 \texttt{FSA}、\texttt{CPA}、\texttt{FPR} 等指标虽然已在服务结果和 README 中预留阈值，但仓库里仍主要是 benchmark hook 和阶段性占位，尚未形成大规模正式实验表。因此，比赛版最合适的表述仍然是“主链已跑通、指标框架已明确、正式实验正在补齐”，而不是把这些阈值写成已经全面达成。

\subsection{短期迭代重点}

比赛版建议优先完成以下工作：
\begin{enumerate}[leftmargin=2em]
  \item 稳定个人决策 DAG 可视化；
  \item 强化 Branch A / Branch B 报告与可逆性曲线；
  \item 补充若干标准化场景样例与对应截图；
  \item 用 benchmark 和测试结果支持 T0 与端到端性能表述；
  \item 统一首页、标题和答辩中的主叙事口径。
\end{enumerate}

\subsection{中长期扩展方向}

在比赛版核心闭环稳定后，可进一步扩展：
\begin{enumerate}[leftmargin=2em]
  \item 更丰富的个体风险因子建模与动态更新；
  \item 更强的多模态输入支持，如截图 OCR、语音转写和证据片段提取；
  \item 更完整的知识库与标注体系；
  \item 更深层的多智能体社会环境建模，用于复杂关系链和长期演化场景；
  \item 更正式的评测基准与专家打分协议。
\end{enumerate}

\section{汇报建议}

\subsection{建议的开场句}

\begin{quote}
\textbf{本项目想解决的不是“这条消息像不像诈骗”，而是“这个人在当前高压情境下会不会马上做出高风险动作，以及怎样在动作发生前改变这条决策路径”。}
\end{quote}

\subsection{建议的主结构}

答辩或汇报时建议采用如下顺序：
\begin{enumerate}[leftmargin=2em]
  \item 为什么基于 MiroFish 可以构建个人决策图谱；
  \item 同类产品与开源方案做到了哪里；
  \item 为什么系统最终采用 T0 / 风险向量 / RAG / Fork / Reporter 这一设计；
  \item 系统结构、五屏工作台与主链工作流；
  \item 当前工程实现、最小 Demo 与 benchmark 基础；
  \item 创新点成立条件、风险治理与后续计划。
\end{enumerate}

\subsection{建议的收尾句}

\begin{quote}
\textbf{本项目的目标不是替代所有反诈产品，而是在损失真正发生之前，为个体提供一张可解释的决策图谱，并通过反事实推演把“风险”转化为“可以被改变的路径”。}
\end{quote}

\section{结论}

\ProjectName{} 的核心价值，不在于把 MiroFish 原有能力机械搬到安全场景，而在于基于其复杂情境建模底座，构造了一条更聚焦、更可解释、也更适合信安竞赛的问题主线：\CorePosition{}。在这一主线下，MiroFish 不再作为宏大世界模拟的展示外壳，而成为构建局部状态、关键分叉和路径比较的运行时基础。

通过对市场竞品与开源系统的调研可以看到，内容识别、企业训练、SOC triage 和社工仿真都已有成熟路径，但围绕“个人关键动作前的反事实决策辅助”仍存在清晰空位。结合双系统理论、用户情境研究、社会工程说服机制研究和反事实解释研究，可以进一步说明：当前项目采用 T0 / 风险向量 / RAG / Fork / Reporter 的组合，并不是功能堆叠，而是对问题本质、竞品边界与理论基础共同约束后的结果。

从工程上看，系统已经具备前后端骨架、主链服务、最小 Demo、五屏工作台、测试与 benchmark 钩子，说明其并非停留在概念层。下一阶段最关键的任务，是把“关键节点、双分支、可逆性、处方”这一主线进一步可视化、量化和标准化。只要这一链路稳定成立，\ProjectName{} 就能够以一套独立、完整、可信的形式支撑竞赛提交、现场答辩和后续深化研发。

\end{document}
